% sections/risks.tex -- risks, rollout phases, and conclusion
\section{Risks}
\label{sec:risks}

\paragraph{Wire compatibility breaks.}
Every Tier-1 migration changes the byte layout of a consensus-critical
or RPC-public surface. The Quasar activation at unix~$1766708400$
(2025-12-25T16:20:00-08:00) is the single coordinated cut-over. Two
mitigations:
\begin{itemize}[leftmargin=*]
  \item Pre-activation P/X blocks remain decodable through the
    \texttt{v0} legacy decoder in
    \texttt{node/vms/platformvm/block/v0/block.go}~--- this is the
    \emph{only} retained legacy code path and it is read-only.
  \item Post-activation, the canonical
    \texttt{luxfi/genesis/configs/\{brand\}-\{net\}/genesis.json}
    embeds the ZAP-native chain configuration at block 0 (memory
    \texttt{canonical\_genesis\_only}). No RLP imports, no fallback.
\end{itemize}

\paragraph{External boundary regressions.}
Tier-4 surfaces (Stripe, Plivo, OpenAI, Binance) keep their vendor
wire. The risk is over-zealous ZAP-ification of an outbound edge,
breaking a third-party integration. Mitigation: the ZAP boundary
terminates at the last Hanzo-owned hop. Inside the gateway, internal
Hanzo$\to$Hanzo calls are ZAP-tunnelled; outbound to external vendors
keeps the vendor format.

\paragraph{Operational visibility.}
JSON-RPC traffic is grep-able in logs and tcpdump. ZAP frames are
opaque binary. Mitigation: the \texttt{luxfi/zap} package ships a
canonical \texttt{zap dump} decoder; gateway access logs include a
ZAP-frame summary line so on-call retains diagnostic parity. The
schema registry lives in-tree (one struct = one schema = one wire),
so frame interpretation is deterministic without out-of-band
schema fetch.

\paragraph{Tier-3 duplicate code.}
The \texttt{lux/proto/} subtree shadows \texttt{lux/protocol/} for 31
files. Both currently import \texttt{luxfi/codec}. The risk is that
migrating one and not the other leaves a half-migrated workspace.
Mitigation: \texttt{lux/proto/} is deleted as Phase~D, not migrated.

\paragraph{TS encoder version skew.}
\texttt{hanzo/zap-es} is published to npm and used by browser-side
clients (extension, chat UI). Wire schema bumps must update both
\texttt{luxfi/zap} (Go) and \texttt{zap-es} (TS) in lockstep.
Mitigation: \texttt{zap-es} bumps follow patch-only semver
(\texttt{1.x.y} only, per workspace rule); breaking changes require a
coordinated major bump that is currently blocked by the
``forwards-perfection only'' rule.

\section{Phased rollout summary}
\label{sec:rollout}

\begin{center}
\begin{tabular}{lll}
\toprule
\textbf{Phase} & \textbf{Surfaces} & \textbf{Dependency} \\
\midrule
A & utxo + protocol/p + protocol/x + node    & none (LP-023 measured) \\
B & agent + mcp + jin tensor wire             & none (parallel with A) \\
C & zen/gateway internal tunnels              & none (independent) \\
D & delete lux/proto + verify zero codec imports & after A \\
\bottomrule
\end{tabular}
\end{center}

Phases A, B, and C can proceed in parallel; only Phase~D blocks on
Phase~A.

\section{Conclusion}
\label{sec:conclusion}

ZAP is the canonical wire for everything Hanzo, Lux, and Liquidity
\emph{own}. Five tiers classify every observed wire surface; Tier 5
(eighteen services) is already done, Tier 4 (three external boundaries)
is deliberately left alone, and Tier 1 (six urgent services) is the
focus of this paper. Measured ratios on adjacent already-migrated
surfaces project a $7.4$--$8.6\times$ compounded speedup and a
$24\times$ allocation reduction on the AI/ML/blockchain hot path. The
proposal does not introduce backwards-compat shims, does not require
re-architecting any service, and composes orthogonally on top of the
Quasar activation cutover already scheduled for unix~$1766708400$.

One wire format, one composition root per surface, one schema registry,
one canonical reference implementation. ZAP everywhere is not a
performance optimisation; it is a decomposition step that takes wire
heterogeneity off the maintenance budget and pays the cost back as
allocation, GC, and tail-latency reduction at every layer of the stack.
